\chapter{}

Quando Anna aveva avuto le convulsioni, restava tutto il giorno stesa nella poltrona; dentro la trattoria. Il suo volto doventava bianco; e Rebecca, assistendola, le slacciava il busto. Ma siccome i cuochi e i camerieri avevano sempre qualche cosa da chiederle, ella riapriva gli occhi, guardava fisso; e poi, scuotendosi tutta, rispondeva. Perchè il marito non s’inquietasse di più, non voleva andare a letto. Ma in quei momenti sentiva una grande angoscia, perchè era incapace di badare a Pietro.

Le sembrava di non appartenere più alla vita, di non avere mai fatto niente per lui. E allora quella specie di quiete, che le dava l’agiatezza, era sempre sciupata dal ricordo della sua miseria. Ella diceva:

— È impossibile esser contenti come vorremmo!

E la stanchezza di esser vissuta era così amara che aveva paura di non sentirsi più buona. Il sentimento della morte le era sempre presente, e non le bastava credere in Dio.

Ella si metteva a guardare Pietro con questo sentimento, e ne provava uno sconforto che le faceva perfino paura.

I suoi nervi scossi dalla convulsione le prolungavano un senso indefinibile di dolore desolato; perchè era avvezza a dover guarire da sè, senza sentire mai che gli altri potevano farle qualche cosa.

Ma sperava di guarire, non perchè credesse al medico, ma perchè aveva Pietro.

Ella non gli sapeva parlare; capiva ch'egli cresceva senza che riuscisse a farselo proprio suo, a dirgli almeno una di quelle parole che avrebbero dovuto consolarla. Anche quando l'aveva vicino, restavano come due che avessero l'impossibilità d'intendersi.

Pietro evitava sempre di farle sentire che le voleva bene, per paura di doventare troppo obbediente; ed ella si disperava troppo e senza ragione di qualche sua scappata. E perciò Pietro temeva quando gli aveva tante cure. Mentre ella, non avendogliele potute fare, cercava un'altra volta d'imporgliele.

— Tu non rispetti la mamma!

Egli, allora, si esasperava; svignandosela senza nè meno ascoltarla.

Anna ci piangeva, dicendolo a Rebecca; che le domandava, con un mezzo sorriso:

— Ma perchè se la prende così?

E siccome glielo aveva allattato e desiderava che fosse affezionato, anche a lei, ci sentiva quasi piacere. Ma Anna, mai accortasi di questo, rispondeva:

— Non lo devi scusare tu!

— Io?

E Rebecca era per offendersi.

Quando poi Pietro la vedeva piangere, credendo che fosse cattiva, gli veniva voglia di far peggio.

Anna consigliava Rebecca e Masa come dovevano educare Ghìsola: era, però, una bontà da padrona; perchè così anche lei dipendeva di più dalla sua volontà. Benchè, le avesse da vero certi riguardi delicati, come quando diceva a Masa che non la facesse lavorare troppo; e come quando, per capo d'anno, pensava sempre a regalarle un vestituccio nuovo, comprato su quei barroccini di merciai che si fermavano all'uscio della trattoria.

Ghìsola, allora, le portava un mazzo di fiori, che, per averli, andava magari a rubare; e le faceva gli augurii.
